%%%%%%%%%%%%%%%%%%%%%%%%%%%%%%%%%%%%%%%%%%%%%%%%%%%%%%%%%%%%%%%%%%%%%%
% Lớp 10 | Phương pháp tọa độ trong mặt phẳng (𝑂𝑥𝑦) | Tọa độ của véctơ
% Dạng 0H9?1-4: Sự cùng phương của 2 véctơ và ứng dụng
% Số câu: 12  (sinh tự động bởi tach_cau.py; từ nay thầy sửa trực tiếp file này)
%%%%%%%%%%%%%%%%%%%%%%%%%%%%%%%%%%%%%%%%%%%%%%%%%%%%%%%%%%%%%%%%%%%%%%

%%%%%%%%%%%%%%%% TRẮC NGHIỆM 4 PHƯƠNG ÁN %%%%%%%%%%%%%%%%

% nguồn: MaHoaNH_TracNghiem_Lop10
\begin{ex}%[0H9H1-4]%[ma:b18ad5a8]
	Giá trị của $m$ sao cho $\overrightarrow{a}=(2m-1;3m)$ cùng phương  $\overrightarrow{b}=(1;1)$ là
	 \choice
		 {$-2$}
		{$1$} 
		{$-3$} 
		{\True  $-1$}
	 \loigiai{
	 $\overrightarrow{a}=(2m-1;3m)$ cùng phương $\overrightarrow{b}=(1;1)$ khi và chỉ khi $$\dfrac{2m-1}{1}=\dfrac{3m}{1}\Leftrightarrow 2m-1=3m\Leftrightarrow m=-1.$$}
\end{ex}

% nguồn: MaHoaNH_TracNghiem_Lop10
\begin{ex}%[0H9H1-4]%[ma:fa5dd40f]
	Trong mặt phẳng với hệ trục tọa độ $Oxy$, cho $\overrightarrow{a}=(1;-2)$. Trong các vectơ dưới đây, vectơ nào cùng phương  với $\overrightarrow{a}$.
	\choice 
		 {$\overrightarrow{b}=(1; 2)$}
		{$\overrightarrow{c}=(-1;-2)$} 
		 {\True $\overrightarrow{d}=(-2; 4)$}
		 {$\overrightarrow{e}=(2; 1)$}
	\loigiai{ 
		Ta có $\overrightarrow{d}=(-2; 4)=-2(1;-2)=-2 \overrightarrow{a}$. Do đó, $\overrightarrow{d}$ cùng phương với $\overrightarrow{a}$. 
	} 
\end{ex}

% nguồn: MaHoaNH_TracNghiem_Lop10
\begin{ex}%[0H9V1-4]%[ma:fb24e11a]
	Trong hệ tọa độ $Oxy$, cho ba điểm $A(1;1)$, $B(3;2)$, $C(6;5)$. Tìm tọa độ điểm $D$ để tứ giác $ABCD$  là hình bình hành.
	 \choice
		{$D(4;3)$} 
		{$D(3;4)$} 
		{\True  $D(4;4)$}
		 {$D(8;6)$}
	 \loigiai{
		Gọi  $D(x;y)$. Ta có $\heva{&\overrightarrow{AB}=(2;1) \\&\overrightarrow{DC}=(6-x;5-y).}$\\
	Tứ giác $ABCD$ là hình bình hành $\Leftrightarrow \overrightarrow{AB}=\overrightarrow{DC}$  $\Leftrightarrow \heva{&2=6-x \\&1=5-y}\Leftrightarrow \heva{&x=4 \\&y=4}\Rightarrow D(4;4)$.}
\end{ex}

% nguồn: MaHoaNH_TracNghiem_Lop10
\begin{ex}%[0H9V1-4]%[ma:05e147dd]
	Cho $A(1;2)$, $B(-2;6)$. Tìm tọa độ điểm $M$ thuộc trục $Oy$ sao cho ba điểm $A$, $B$, $M$ thẳng hàng? 
	\choice 
		 {$M(0;3)$}
		 {\True $M\left(0;\dfrac{10}{3}\right)$}
		 {$M\left(\dfrac{5}{2};0\right)$}
		 {$M\left(0;\dfrac{5}{2}\right)$}
	\loigiai{ 
		 Vì $M$ thuộc trục $Oy$ nên $M(0;y)$.
		 Suy ra $\overrightarrow{AB}=(-3;4)$, $\overrightarrow{AM}=(-1;y-2)$.\\
		Để ba  điểm $A$, $B$, $M$ thẳng hàng thì $\dfrac{-1}{-3}=\dfrac{y-2}{4}\Leftrightarrow 4=3y-6$ $\Leftrightarrow y=\dfrac{10}{3}$.\\
	Vậy $M\left(0;\dfrac{10}{3}\right)$.} 
\end{ex}

% nguồn: MaHoaNH_TracNghiem_Lop10
\begin{ex}%[0H9C1-4]%[ma:9a59e5ec]
	Biết rằng hai vectơ $\overrightarrow{a}$ và $\overrightarrow{b}$ không cùng phương nhưng  hai vectơ $2\overrightarrow{a}-3\overrightarrow{b}$ và $\overrightarrow{a}+(x-1)\overrightarrow{b}$ cùng phương. Khi đó giá trị của $x$ là
	 \choice
		 {$\dfrac{1}{2}$}
		{$-\dfrac{3}{2}$} 
		 {\True $-\dfrac{1}{2}$}
		{$\dfrac{3}{2}$} 
	\loigiai{ 
		Do hai vectơ $2\overrightarrow{a}-3\overrightarrow{b}$ và $\overrightarrow{a}+(x-1)\overrightarrow{b}$  cùng phương.\\
		 Suy ra $2\overrightarrow{a}-3\overrightarrow{b}=k\left[\overrightarrow{a}+(x-1)\overrightarrow{b}\right]\,\left(k\ne 0, k\in \mathbb{R}\right)$\\
		$\Rightarrow 2\overrightarrow{a}-3\overrightarrow{b}=k\overrightarrow{a}+k(x-1)\overrightarrow{b}\Rightarrow  (k-2)\overrightarrow{a}+\left[k(x-1)+3\right]\overrightarrow{b}=\overrightarrow{0}$.\hfill$(1)$\\
		Theo đầu bài hai  vectơ $\overrightarrow{a}$ và $\overrightarrow{b}$ không cùng phương.
		 $$(1)\Rightarrow \heva{&k=2 \\&k(x-1)=-3}\Leftrightarrow \heva{&k=2 \\&x-1=-\dfrac{3}{2}}\Leftrightarrow \heva{&k=2 \\&x=-\dfrac{1}{2}.}$$
	Vậy $x=-\dfrac{1}{2}$.} 
\end{ex}


%%%%%%%%%%%%%%%% ĐÚNG/SAI %%%%%%%%%%%%%%%%

% nguồn: MaHoaNH_DungSai_Lop10
\begin{ex}%[0H9N1-4]%[ma:e7bb6f4b]
	Trong mặt phẳng tọa độ $Oxy$. $\overrightarrow{m}=(1;1)$ cùng phương với vec-tơ nào dưới đây 
	\choiceTF 
		 {\True $\overrightarrow{OA}=(2;2)$}
		{$\overrightarrow{OB}=(0;1)$} 
		{\True  $\overrightarrow{OC}=(-3;-3)$}
		 {$\overrightarrow{OD}=(-1;0)$}
	\loigiai{ 
		\begin{itemchoice}
			\itemch $\overrightarrow{m}=(1;1)$ cùng  phương với $\overrightarrow{OA}=(2;2)$.
			 \itemch $\overrightarrow{m}=(1;1)$ không cùng phương với $\overrightarrow{OB}=(0;1)$.
			\itemch $\overrightarrow{m}=(1;1)$ cùng phương với $\overrightarrow{OC}=(-3;-3)$. 
			\itemch $\overrightarrow{m}=(1;1)$ không cùng phương với  $\overrightarrow{OD}=(-1;0)$.
		\end{itemchoice}
	 }
\end{ex}

% nguồn: MaHoaNH_DungSai_Lop10
\begin{ex}%[0H9N1-4]%[ma:0d1a86e4]
	Trong mặt phẳng tọa độ $Oxy$. Tồn tại số  thực $m$ để $\overrightarrow{OM}=(m;1)$ cùng phương với
	 \choiceTF
		{\True $\overrightarrow{OA}=(2;2)$} 
		{$\overrightarrow{OB}=(0;1)$} 
		 {\True $\overrightarrow{OC}=(-3;-3)$}
		{$\overrightarrow{OD}=(-1;0)$} 
	\loigiai{ 
		\begin{itemchoice}
			 \itemch $m=1$ thì $\overrightarrow{OM}=(1;1)$ cùng phương với $\overrightarrow{OA}=(2;2)$.
			\itemch $m=0$ thì  $\overrightarrow{OM}=(0;1)$ cùng phương với $\overrightarrow{OB}=(0;1)$.
			\itemch $m=1$ thì  $\overrightarrow{OM}=(1;1)$ cùng phương với $\overrightarrow{OC}=(-3;-3)$.
			\itemch Không  tồn tại $m$ để $\overrightarrow{OM}=(m;1)$ cùng phương với $\overrightarrow{OD}=(-1;0)$.
		\end{itemchoice}
	 }
\end{ex}

% nguồn: MaHoaNH_DungSai_Lop10
\begin{ex}%[0H9V1-4]%[ma:59a78bd2]
	Trong mặt phẳng tọa độ $(Oxy)$, cho $3$ điểm $A(-2;-1)$, $B(1;3)$, $C(2;-3)$. Xét tính Đúng/Sai của  các mệnh đề sau
	\choiceTF 
		{Độ dài  đoạn thẳng $AB=25$}
		{\True $A$, $B$, $C$ là  ba đỉnh của một tam giác}
		{Điểm $I(0;-2)$ là trung điểm của $AB$} 
		{Điểm $M$  thuộc $Ox$ sao cho $AM+BM$ bé nhất có hoành độ bằng $\dfrac{5}{4}$}
	 \loigiai{
		\begin{center}
			\begin{tikzpicture}[line join = round,line cap = round,>=stealth,font=\footnotesize,scale=0.8]
					\draw[thick,->] (-4,0)--(5,0) node  [below]{$x$};
					 \draw[thick,->] (0,-4)--(0,4) node [below left]{$y$};
					 \foreach \i in {-3,...,3}{\draw[fill=black](\i,0) circle (1.2pt) node[below left]{$\i$};}
					\foreach \i in {-3,...,3}{\draw[fill=black](0,\i) circle (1.2pt) node[below left]{$\i$};} 
					\path 
						(-2,-1) coordinate  (A)
						(1,3) coordinate  (B)
						 (2,-3) coordinate (C)
						($(A)!0.5!(B)$)  coordinate (I);
					\draw[dashed] (-2,0)--(A)--(0,-1) 
						 (1,0)--(B)--(0,3)
						 (2,0)--(C)--(0,-3);
					 \draw (A)--(I) node[midway,sloped]{\tiny |}--(B) node[midway,sloped]{\tiny |}--(C)--cycle;
					\foreach  \x/\g in {A/180,B/45,C/0,I/135}
						{\draw [fill=black](\x) circle  (1.2pt) +(\g:3mm) node{$\x$};}
				\end{tikzpicture}
		\end{center}
		\begin{itemchoice}
			\itemch Độ dài  đoạn thẳng $AB=\sqrt{(1+2)^2+(3+1)^2}=5$.\\
			Mệnh đề  sai.
			\itemch Ta có $\overrightarrow{AB}=(3;4)$, $\overrightarrow{AC}=(4;-2);$ vì $\dfrac{3}{4}\ne \dfrac{4}{-2}\Rightarrow \overrightarrow{AB}$ và  $\overrightarrow{AC}$ không cùng phương.\\
			Mệnh  đề đúng.
			\itemch Trung điểm của $AB$ là  $I \left( \dfrac{-1}{2};1\right)$.\\
			Mệnh  đề sai.
			\itemch Do $y_A\cdot y_B=(-1)\cdot 3<0$ nên hai  điểm $A$, $B$ nằm khác phía so với $Ox$.
			\begin{center}
				\begin{tikzpicture}[line join = round,line cap = round,>=stealth,font=\footnotesize,scale=0.8]
						\draw[thick,->]  (-4,0)--(5,0) node [below]{$x$};
						\draw[thick,->] (0,-2)--(0,4)  node [below left]{$y$};
						 \foreach \i in {-3,...,3}{\draw[fill=black](\i,0) circle (1.2pt) node[below left]{$\i$};}
						\foreach \i in {-1,...,3}{\draw[fill=black](0,\i) circle (1.2pt) node[below left]{$\i$};} 
						\path 
						(-2,-1) coordinate (A) 
						(1,3) coordinate (B) 
						(2,-3)  coordinate (C)
						(1.5,0) coordinate (M) 
						($(A)!0.5!(B)$) coordinate (I); 
						\draw[dashed]  (-2,0)--(A)--(0,-1)
						 (1,0)--(B)--(0,3);
						\draw (A)--(I) node[midway,sloped]{\tiny |}--(B) node[midway,sloped]{\tiny  |}
						 (A)--(M)--(B);
						\foreach  \x/\g in {A/180,B/45,I/135,M/45}
						{\draw [fill=black](\x) circle  (1.2pt) +(\g:3mm) node{$\x$};}
					\end{tikzpicture}
			\end{center}
			Do đó, $M\in Ox$ để $AM+BM$ nhỏ nhất $\Rightarrow A,B,M$ thẳng hàng  $\Rightarrow \overrightarrow{AB}$ và $\overrightarrow{AM}$ cùng phương.\\
			Gọi $M(x;0)\in Ox\Rightarrow \overrightarrow{AM}=(x+2;1)$ và $\overrightarrow{AB}=(3;4)$.\\ 
			Ta có  $\overrightarrow{AM},\overrightarrow{AB}$ cùng phương $\Rightarrow \dfrac{x+2}{3}=\dfrac{1}{4}\Rightarrow x=-\dfrac{5}{4}$.\\
			Vậy  $M \left(-\dfrac{5}{4};0 \right)$.\\
			Mệnh đề sai. 
		\end{itemchoice}
	} 
\end{ex}


%%%%%%%%%%%%%%%% TRẢ LỜI NGẮN %%%%%%%%%%%%%%%%

% nguồn: MaHoaNH_TraLoiNgan_Lop10
\begin{bt}%[0H9H1-4]%[ma:9940787e]
	Trong mặt phẳng toạ độ $Oxy$, cho hai vectơ $\overrightarrow{a}=(3 m; 4 m-1)$ và $\overrightarrow{b}=(\sqrt{2}; \sqrt{2})$ (với $m$ là tham số). Tìm $m$ để hai vectơ  $\overrightarrow{a}$ và $\overrightarrow{b}$ cùng phương.
	\shortans{1} 
	 \loigiai{
		Ta có $\overrightarrow{a}$  và $\overrightarrow{b}$ cùng phương khi và chỉ khi $3m=4m-1\Leftrightarrow m=1$.
	} 
\end{bt}

% nguồn: MaHoaNH_TraLoiNgan_Lop10
\begin{bt}%[0H9H1-4]%[ma:30d6731f]
	Cho $\overrightarrow{u}=(2 x-1; 3)$, $\overrightarrow{v}=(1; x+2)$. Có hai giá trị của $x$ để $\overrightarrow{u}$  cùng phương với $\overrightarrow{v}$. Tính tích hai giá trị đó.
	 \shortans{-2,5}
	 \loigiai{
		Với $x=-2$: Ta có $\overrightarrow{u}=(-5; 3); \overrightarrow{v}=(1;  0)$.\\
		Vì  $\dfrac{1}{-5} \neq \dfrac{0}{3}$ nên hai vectơ $\overrightarrow{u}; \overrightarrow{v}$ không cùng phương.\\
		Với $x \neq-2$: Ta có $\overrightarrow{u}; \overrightarrow{v}$ cùng phương khi và chỉ khi $$\dfrac{2 x-1}{1}=\dfrac{3}{x+2}\Leftrightarrow(2 x-1)(x+2)=3 \Leftrightarrow 2 x^2+3 x-5=0  \Leftrightarrow\hoac{&x=1 \\& x=-\dfrac{5}{2}.}$$
		 Vậy tích của chúng là $1 \cdot\left(-\dfrac{5}{2}\right)=-2{,}5$.
	} 
\end{bt}

% nguồn: MaHoaNH_TraLoiNgan_Lop10
\begin{bt}%[0H9H1-4]%[ma:afe03bed]
	Cho $\overrightarrow{a}=(1;3), \overrightarrow{b}=(6;-2)$. Một $\overrightarrow{d}$ thỏa mãn $\overrightarrow{a} \perp \overrightarrow{d}$  và $\overrightarrow{b} \cdot \overrightarrow{d}=20$ thì giá trị biểu thức $x+y$ bằng
	 \shortans{2}
	 \loigiai{
		Ta có $\begin{aligned}
		&\heva{&\overrightarrow{a} \perp \overrightarrow{d} \\&\overrightarrow{b} \cdot \overrightarrow{d}=20}  \Leftrightarrow\heva{&\overrightarrow{a} \cdot \overrightarrow{d}=0 \\&\overrightarrow{b} \cdot \overrightarrow{d}=20}\\
		&\Leftrightarrow\heva{&x+3 y=0 \\&6  x-2 y=20} \Leftrightarrow\heva{&x=3 \\&y=-1}\Rightarrow \overrightarrow{d}=(3;-1).
	\end{aligned}$\\
	Vậy $x+y=3+(-1)=2$. 
} 
\end{bt}


%%%%%%%%%%%%%%%% TỰ LUẬN %%%%%%%%%%%%%%%%

% nguồn: MaHoaNH_TuLuan_Lop10
\begin{bt}%[0H9V1-4]%[ma:b9334d48]
	Cho các $\overrightarrow{a}=(1;-2)$,  $\overrightarrow{b}=(-2;-6)$, $\overrightarrow{c}=(m+n;-m-4 n)$.
	Tồn tại hai số thực  $m$, $n$ sao cho $\overrightarrow{c}$ cùng phương $\overrightarrow{a}$ và $|\overrightarrow{c}|=3 \sqrt{5}$. Tìm giá trị của biểu thức $|m+n|$.
	 \loigiai{
		$ \overrightarrow{c} $ cùng phương $ \overrightarrow{a} $ và $ \left|\overrightarrow{c}\right|=3 \sqrt{5} $  nên
		\begin{align*}
			\heva{&\dfrac{m+n}{1}=\dfrac{-m-4 n}{-2} \\ &\sqrt{(m+n)^2+(-m-4  n)^2}=3 \sqrt{5}}
			 &\Leftrightarrow
			\heva{&- 2 m - 2 n = - m  - 4 n \\ &( m + n) ^ { 2 } + ( m + 4 n) ^ { 2 } = 4 5} \\
			&\Leftrightarrow 
			\heva{&m = 2 n \\ &( 3 n) ^ { 2 } + ( 6 n) ^ { 2  } = 4 5} \\
			 &\Leftrightarrow
			\heva{&m=2 n \\ &4 5 n ^ { 2 } = 4 5}  \\
			&\Leftrightarrow 
			\heva{&m=2 \\ &n=1}\ \text{hoặc}\ \heva{&m=-2  \\ &n=-1.}
		\end{align*}
		Trong cả  hai trường hợp, ta đều thấy $ |m+n| = 3 $.\\
		Vậy  $|m+n|=3$.
	} 
\end{bt}
